\section{Introduction}
\label{sec:introduction}

Software engineering is a naturally complicated discipline, traditionally
ruled by a sequential, interconnected series of practices that include
requirements analysis, design, implementation, and testing. The Model-View-Controller (MVC) architecture, first introduced 
by Reenskaug~\citep{Reenskaug1979} and later formalized as a 
design pattern~\citep{Gamma1994}, remains one of the most 
widely adopted structural paradigms in software development.
By enforcing a clear separation of concerns across data management, user
interface, and control logic, MVC provides a principled foundation for
building maintainable, testable, and scalable systems. Even in the era
of AI-assisted development, this separation remains relevant; as code
generation accelerates, preserving architectural discipline becomes
increasingly critical.

However, the rapid integration of Large Language Models (LLMs) into the
development toolchain has been reported to erode this
discipline~\citep{Slater2025, Waseem2025, Ehsani2026}. AI-powered coding assistants
such as Cursor~\citep{Cursor2025}, GitHub Copilot~\citep{GitHubCopilot2025}, 
and OpenAI Codex~\citep{OpenAICodex2025} have significantly
accelerated implementation tasks~\citep{Hou2024}. However, this
speed may come at the expense of macro-level architectural
integrity, the structural coherence that governs inter-layer boundaries
and system-wide design consistency~\citep{Whiting2020}. The speed at which LLMs translate
complex business logic into functional code, while effective in isolated
benchmarks, deserves careful evaluation when applied to real-world,
architecture-sensitive projects.

This pattern is increasingly described as ``vibe
coding''~\citep{Fawzy2025}, an approach where developers rely heavily
on LLM outputs, generating code based on momentary decisions
rather than adhering to a defined architecture. The underlying danger is
that the generated code may function correctly in isolation, a function
may yield the correct output, a class may behave as expected, yet when
these components are integrated, the resulting structure often
contradicts the project's architectural principles and silently violates
inter-layer boundaries. LLM-generated code has been found to contain
structural and quality problems, such as poor exception handling,
undefined variables, and weak adherence to design rules, that can lead
to architectural drift in layered systems when combined without careful
review~\citep{Siddiq2024, Fawzy2025, Slater2025}. When developers
commit such violations to the codebase without adequate review,
technical debt accumulates without control.

Over time, this accumulation shows itself as Architecture
Erosion~\citep{Li2022}, defined as the progressive accumulation of
deviations between the intended architecture and the implemented
codebase. Architecture erosion goes beyond structural degradation: it restricts system
scalability and complicates the integration of new
requirements~\citep{Li2022}. Although various enforcement and
remediation approaches have been proposed to detect and address
architecture erosion~\citep{Li2022}, it remains a major unresolved
challenge in software engineering. As this accumulation stems from
architecturally non-compliant code introduced during development~\citep{Li2022},
we reason that mitigating it in AI-assisted workflows requires enforcement
mechanisms acting at the point of code generation, before such violations
can accumulate into long-term drift. Addressing this upstream,
point-of-generation problem, rather than measuring erosion's longitudinal
accumulation, is therefore the specific focus of the present study.

Existing tools have failed to mitigate this problem fully. Traditional
static analysis tools, such as SonarQube~\citep{SonarQube} and
Checkstyle~\citep{Checkstyle}, excel at detecting low-level syntactic
and stylistic anomalies; however, they largely overlook high-level
architectural deviations~\citep{Fontana2017}, such as inter-layer
responsibility leaks or MVC isolation violations. Furthermore, most current
architectural compliance tools require substantial manual preprocessing
to map the system to its intended architecture. This well-documented
barrier often leads practitioners to rely on informal, error-prone
inspection methods~\citep{Ali2018}. Another critical gap is the absence
of a structured, end-to-end traceability chain between Software
Requirements Specifications (SRS) and the generated code. While
existing traceability methods predominantly rely on Information
Retrieval (IR) techniques, they continue to exhibit significant
limitations in terms of full automation and seamless integration into
modern development workflows~\citep{Koboyatshwene2025}.

We propose the \textit{MVC Test Orchestrator} to address this gap, a
multi-agent system designed to preserve and enforce MVC architectural
integrity across the AI-assisted development workflow and built on an
``Architecture-First'' philosophy, wherein all design and code
generation decisions are anchored to and constrained by the SRS before
any implementation step begins. Concretely, the framework: (i)~extracts
MVC components from unstructured SRS text via RAG-grounded multi-agent
orchestration (\texttt{/extract}); (ii)~generates rule-based project
scaffolds that enforce MVC layer boundaries at the file-system level
(\texttt{/scaffold}); (iii)~detects inter-layer architectural violations
through AST-based static analysis (\texttt{/audit}); and
(iv)~autonomously remediates identified violations within the boundaries
of the established Architectural Map (\texttt{/fix}). The system
establishes the SRS document as the single source of truth, governing
this entire workflow from requirements analysis through autonomous
remediation. Rather than functioning as a fully autonomous black box,
the system employs a Human-in-the-Loop (HITL) architecture that
preserves developer oversight and strategic authority at every stage,
acknowledging that LLM-based agents can approximate but not fully
guarantee rule-bound behavior. Underlying these mechanisms is what we
term a Safety-Layered Architecture: rule-based components (scaffolding,
AST-based auditing) are interleaved with the LLM's generative steps,
aiming to bound their stochastic variability within verifiable
architectural constraints.

\smallskip\noindent\textbf{Scope.} The prototype is intended to support MVC structural conformance in newly generated Python artifacts. The study does not observe the longitudinal evolution of a production codebase and therefore does not directly measure architecture erosion, architectural drift over time, or technical-debt accumulation.

This paper makes three contributions. First, we present a multi-agent
architecture that automates the complete MVC lifecycle, from SRS
ingestion to self-healing remediation. This architecture combines
RAG-based semantic grounding with rule-based scaffolding and AST-driven
auditing under a HITL oversight protocol. Second, we introduce an
eight-metric evaluation framework spanning both LLM-based generative
stages and rule-based pipeline stages. These metrics include Semantic
Consistency (BERTScore), Architectural Component Extraction Score
(ACES), Scaffolder Error Rate (SER), Violation Detection Recall and
Alert Confirmation Precision (VDR/ACP), Repair Success among Detected
Alerts and end-to-end Repair Success Rate (RSDA/RSR\textsubscript{e2e}), and
Consistency Score (CS). Third, we report empirical results across 20
SRS documents of varying domain and complexity. The system achieves an
extraction score of 74.0\%, a violation detection recall of 98.8\%
(96.2\% precision), a remediation success rate of 92.6\% among detected
violations, a 0\% scaffolding error rate across all evaluated
documents, and a 70.6\% consistency score across three independent
iterations per document.

The remainder of this paper is organized as follows: Section~\ref{sec:related_work} reviews
the relevant literature, Section~\ref{sec:methodology} describes the proposed methodology
and system architecture, Section~\ref{sec:experiment} presents the experimental
design---including the five Research Questions (RQ1--RQ5) that
structure the evaluation---and the evaluation metrics, Section~\ref{sec:results}
reports the experimental results, Section~\ref{sec:discussion} discusses the findings,
Section~\ref{sec:limitations} addresses the limitations and threats to validity, Section~\ref{sec:conclusion}
concludes the study, and Section~\ref{sec:future} outlines potential directions for
future research.

\smallskip\noindent\textbf{Extension beyond the preliminary conference paper.} This journal article substantially extends the five-case preliminary study reported by \cite{dim2026mvc}. The evaluation corpus is expanded from 5 to 20 SRS documents, and the architecture is redesigned around three specialized Architect Agents instead of one. This study also adds an autonomous remediation stage that was not present in the conference version, and replaces the earlier AI-as-a-Judge evaluation with a five-expert human panel supported by an Inter-Rater Reliability analysis. The evaluation framework itself grows from five simple metrics to eight, organized around five research questions (RQ1--RQ5) not present in the earlier work. The conference paper is cited wherever the original extraction and scaffolding concepts are described; reused material is limited to the underlying problem framing, the core RAG mechanism, and the public code repository.